% !TeX root = ../main.tex
\section{Style transfer, effect representations and sound matching}
\label{sec:matching}

\subsection{Reference-based style transfer}

Audio-effect \term{style transfer} processes an input so that it sounds like a
reference. DeepAFx-ST predicts EQ and compressor parameters from an input and a
reference, trained self-supervised with an audio loss; numerical gradients (SPSA)
let it drive non-differentiable effects at a cost in
stability~\cite{steinmetz2022style}. Grant trained the same way through legacy
plugins~\cite{grant2023style}. Koo et al.\ learn an effects encoder by
contrastive learning and render input stems in a reference's mixing style; its
embedding distances on Demucs-separated stems stayed close to those on clean
stems~\cite{koo2023music}. Diff-MST predicts a differentiable mixing console for
any number of tracks~\cite{vanka2024diffmst}.

\textbf{ST-ITO} removes the trained controller: CMA-ES searches the parameters of
an arbitrary plugin chain to maximise the cosine similarity of AFx-Rep embeddings
of the processed input and the reference~\cite{steinmetz2024stito}. AFx-Rep is
trained without dry references to recognise which of 63 plugins and 10 presets
produced an output, and beats CLAP and MFCCs on zero-shot style classification
(0.86 against 0.56 for CLAP) \reported. The search uses a population of 64 for at
most 25 steps (about a minute). The authors state that it does not work well for
challenging cases such as \emph{guitar tone matching}. Yu et al.\ add a Gaussian
prior fitted on real presets, turning the search into a maximum-a-posteriori
estimate and reducing parameter error by up to a third~\cite{yu2025improving}.
ITO-Master found that optimising a learned conditioning embedding beats
optimising raw DSP parameters for mastering~\cite{koo2025itomaster}.

\begin{caution}[: an embedding is not an objective]
TAPE predicts editable parameters of its own differentiable chain (EQ,
compressor, waveshaper, a small residual network, reverb) in one pass and runs in
real time~\cite{long2026tape}. Its comparison is the most instructive result for
us: CMA-ES over ST-ITO's chain with the AFx-Rep objective lowered the embedding
distance (0.259 against 0.660 for TAPE) but raised the spectral error to 2.91,
\emph{above} the 2.34 of the unprocessed input, while TAPE reached 1.12
\reported. An embedding that is good for retrieval can be a bad objective to
optimise. The app's search therefore uses a spectral distance.
\end{caution}

\subsection{Effect representations}

Several encoders describe effects independently of content: AFx-Rep, the effects
encoder of Koo et al., Fx-Encoder++, the tone encoder of Chen et
al.~\cite{chen2024zeroshot}, the capture-trained encoder of Open-Amp (91.1\%
effect-class accuracy on EGDB~\cite{wright2025openamp}) and TAPE's.
Fx-Encoder++ extracts an instrument-wise effects embedding from a mixture given
an audio or text query~\cite{yeh2025fxencoder}, and gives \textbf{the only
measurement on separated stems}: effect retrieval R@10 falls from 39.4\% on true
isolated stems to 20.5\% on Demucs-separated stems \reported, attributed to
separation artefacts distorting effect cues. RelFx learns the \emph{relative}
effect between two processed clips, so it needs no dry
reference~\cite{liu2026beyond}. General embeddings (OpenL3, PANNs, CLAP) move
monotonically with effect strength but not linearly, so they are weak parameter
estimators~\cite{deng2025investigating}. TimberAgent retrieves the nearest of
1\,063 guitar-effect presets with Gram matrices of Wav2Vec2
features~\cite{he2026timberagent}.

\subsection{Text and language models}

Text2FX optimises EQ and reverb against a CLAP text embedding~\cite{chu2025text2fx};
LLM2Fx predicts them with language models, better than Text2FX on guitar EQ
(distribution distance 0.15 against 0.42)~\cite{doh2025llm2fx}; FxSearcher uses
Bayesian optimisation over a plugin chain~\cite{ki2025fxsearcher}. The ``Does it
chug?'' dataset collects guitar-tone adjectives on clips that vary EQ and
distortion~\cite{sutar2024chug}. None estimates amp or drive settings; text suits
refinement (``more bite'') better than estimation.

\subsection{Lessons from synthesizer sound matching}

Finding synthesizer settings that reproduce a sound is the same inverse problem
with a longer history. Its findings carry over:
\begin{itemize}
\item \textbf{Loss choice}: training on audio losses alone often fails for
  frequency-like parameters; pre-training on a parameter loss and fine-tuning on
  a spectral loss works best~\cite{masuda2021ddsp,uzrad2024diffmoog}. LFO rates,
  delay times and wah positions are such parameters.
\item \textbf{Optimisers}: on non-convex landscapes CMA-ES beat Adam (loss 2.09
  against 4.60)~\cite{bogdan2026instrumental}; evolution strategies match sounds
  through non-differentiable synthesizers with a CLAP objective~\cite{cherep2024ctag};
  reinforcement learning with rewards from rendered audio improves
  out-of-domain matching~\cite{wu2026ddsynthrl}.
\item \textbf{Ambiguity}: when different settings give the same sound, point
  regression fails and generative models of the parameters do
  better~\cite{hayes2025equivariant}.
\item \textbf{Proxies}: a network that maps presets to a pretrained audio
  embedding makes black-box synthesizers trainable~\cite{combes2025proxies}.
\item \textbf{Domain gap}: every system degrades on real sounds outside its
  training renders~\cite{hayes2024review}.
\end{itemize}

\begin{keyidea}[: the consensus]
The effect and synthesizer literatures agree on a hybrid: predict an initial
setting with a network trained on renders of the target chain, refine it by
gradient-free search through the real renderer, regularise with a prior over
realistic presets, and keep several candidates.
\end{keyidea}
